
\def\PUDcodigo{EME105}

\if\COMPILAR0
    \def\PUDcodacad{04507.6}
\else
    \def\PUDcodacad{04506.6}
\fi

\def\PUDdisciplina{Desenho Técnico}

\def\PUDsemestre{S1}

\def\PUDhoras{80}

%NOVOS ---------------------------------------------------
\def\PUDhorasTeoria{30}
\def\PUDhorasPratica{50}

\if\COMPILAR0
   \def\PUDinterECA{\hyperref[PUD:EME106]{Desenho Auxiliado por Computador (04507.11)} 
   
   \hyperref[PUD:EME107]{Metrologia (04507.16)} 
   
   \hyperref[PUD:EME123]{Manufatura Auxiliada por Computador (04507.46)}
   }
\else
   \def\PUDinterECA{\hyperref[PUD:EME106]{Desenho Auxiliado por Computador (04506.10)} 
   
   \hyperref[PUD:EME107]{Metrologia (04506.14)} 
   
   \hyperref[PUD:EME123]{Manufatura Auxiliada por Computador (04506.42)}
   }
\fi

\def\PUDatuaisECA{---}

\def\PUDrecursos{Projetor multimídia; Quadro branco e pincel; Aulas em laboratório de informática.}

%----------------------------------------------------------

\def\PUDcreditos{4}

\def\PUDprerequisito{---}

\def\PUDpreacad{---}

\def\PUDobjetivos{Ler, interpretar e reproduzir um desenho técnico. Conhecer os materiais e normas utilizadas no desenho técnico. Compreender as vistas ortográficas, cortes e secções de um objeto e sua representação em perspectiva.
%Dimensionar os principais tipos de transportadores industriais.
}

\def\PUDementa{Introdução ao Desenho Técnico;  Aspectos Gerais do Desenho Técnico;  Perspectivas;  Projeções Ortogonais;  Cotagem;  Escalas;  Corte e Seções;  Tolerância Dimensional e Estado de Superfície.}

\def\PUDconteudo{ &  Introdução ao Desenho Técnico: Importância do desenho técnico, Diferenças entre desenho técnico e desenho artístico, Modos de representação do desenho técnico: perspectivas; vistas múltiplas, Principais normas de desenho Técnico, Visão geral de um sistema CAD em desenho técnico.  Figuras Geométricas; 
\\ 
 &  Aspectos Gerais do desenho técnico: Escrita Normalizada.  Tipos de linha.  Folhas para desenho.  Legendas.  Escalas; 
\\ 
 &  Perspectivas: Objetivo do desenho em perspectiva.  Tipos de perspectivas: Isométrica, Cavaleira, Dimétrica e Cônica.  Eixos Isométricos e Métodos de construção da Perspectiva Isométrica; 
\\ 
 &  Projeções Ortogonais: Conceito de projeção.  Método Europeu e o método Americano de projeções.  Classificação das projeções Geométricas Planas.  Representação em múltiplas vistas.  Significados das linhas.  Vistas necessárias e suficientes e escolha das vistas.  Vistas Parciais, deslocadas e interrompidas; 
\\ 
 &  Cotagem: Aspectos gerais da cotagem, Elementos da cotagem, Disposição das cotas nos desenhos, Cotagem dos elementos, Critérios de cotagem, Cotagem de representações especiais, Seleção das cotas; 
\\ 
 &  Escalas: Objetivo do uso de escalas.  Tipos de Escalas: Natural, de redução e de ampliação.  Aplicação de escalas de redução e de ampliação em desenhos de perspectivas e projeções ortogonais; 
\\ 
 &  Cortes e Seções:  Modos de cortar as peças, Cortes por planos paralelos ou concorrentes, Regras gerais em corte, Elementos que não são cortados e representações convencionais, Cortes em desenhos de conjunto de peças, Seções e encurtamento; 
\\ 
 &  Tolerância dimensional e Estados de Superfície: Tolerância dimensional, Sistemas ISO de tolerâncias Lineares, Sistemas ISO de Angulares, Inscrições das tolerâncias nos desenhos, Ajustes, Verificação das tolerâncias, Tolerância dimensional Geral, Tolerância de peças especiais, Estados de Superfícies. }

\def\PUDmetodo{Aulas expositórias que podem ser teóricas e/ou práticas, onde as práticas no laboratório serão marcadas no decorrer da disciplina.}

\def\PUDavalia{A avaliação será feita com aplicação de provas teóricas e/ou práticas, além da possibilidade de inclusão de trabalhos e seminários no decorrer da disciplina.}

\def\PUDbibbasica{ & \href{www.biblioteca.ifce.edu.br/index.asp?codigo_sophia=23833}{Maguire}, D. E.  Desenho Técnico: problemas e soluções gerais de desenho.  Volume Único.  Editora Hemus.  Rio de Janeiro, 2004.  257p.  ISBN: 9788528903966. 
\\ 
 & \href{www.biblioteca.ifce.edu.br/index.asp?codigo_sophia=24198}{Strauhs}, Faimara do Rocio.  Desenho Técnico.  Base Editorial.  Volume Único.  Curitiba, 2010.  112p.  ISBN: 9788579055393. 
\\ 
 & \href{www.biblioteca.ifce.edu.br/index.asp?codigo_sophia=23872}{Silva}, Arlindo;  Ribeiro, Carlos Tavares;  Dias, João;  Sousa, Luís.  Desenho Técnico Moderno.  Volume Único.  Editora LTC.  Rio de Janeiro, 2006.  ISBN: 9788521615224. }

\def\PUDbibcomplementar{ &  \href{www.biblioteca.ifce.edu.br/index.asp?codigo_sophia=40465}{Gerhard}, Pahl.  Projeto na engenharia: fundamentos do desenvolvimento eficaz de produtos, métodos e aplicações.  São Paulo, SP : Blucher, 2013. ISBN: 9788521203636.
%& Carmona, Tadeu.  Gerenciamento e desenho de projetos.  Linux New Media do Brasil.  Volume Único.  São Paulo, 2007.  95p.  ISBN: 9788561024017. 
\\ 
 & \href{www.biblioteca.ifce.edu.br/index.asp?codigo_sophia=23779}{Junghans}, Daniel.  Informática Aplicada ao Desenho Técnico.  Volume Único.  Base Editorial.  Curitiba, 2010.  224p.  ISBN: 9788579055478. 
\\ 
 & \href{www.biblioteca.ifce.edu.br/index.asp?codigo_sophia=24129}{Lima}, Claudia Campos Netto Alves de.  Estudo dirigido de AutoCAD 2007.  4º ed.  São Paulo, SP : Érica, 2011. 300p.  ISBN: 9788536501185.
\\
 & \href{www.biblioteca.ifce.edu.br/index.asp?codigo_sophia=57176}{Ribeiro}, Antonio Clelio; Peres, Mauro Pedro.  Curso de desenho técnico e autocad. E-book.  Pearson.  388p.  ISBN: 9788581430843. 
\\
 & \href{www.biblioteca.ifce.edu.br/index.asp?codigo_sophia=58215}{Silva}, Ailton Santos.  Desenho técnico.  E-book.  Pearson.  136p.  ISBN: 9788543010977. }

\def\PUDautor{Samuel Vieira Dias}

\def\PUDdata{2013-04-23}

\def\PUDrevisao{3}

\def\PUDrevisor{José Ciro dos Santos}

\def\PUDdatarevisao{2019-05-15}

\PUDcorpo
